\lecture{16}
\title{Speech}

\begin{document}

\begin{remark}
    During an actual class of 6.300,
    this lecture features videos and demonstrations
    that I do not have access to while writing these notes.
    Thus, these notes may exclude some 6.300 content.
\end{remark}

\subsection{Lecture: Speech as a System}

Speech can be modeled as the output of an LTI system.

\begin{center}
    \frame{\includegraphics[width=0.2\textwidth]{lecture-16/speech-system}}
\end{center}

The input $x(t)$ is roughly a periodic impulse train
generated by the vocal cords.
The behavior of the filter $H(\omega)$
depends on the shape of the speaker's throat and nasal cavities.

\subsection{Lecture: Vowels and Formants}

Different vowels have different harmonic content.

\begin{center}
    \frame{\includegraphics[width=0.4\textwidth]{lecture-16/vowel-harmonics}}
\end{center}

Recall that the harmonic content $Y(\omega)$ pictured above
is just the product $Y(\omega) = X(\omega) \cdot H(\omega)$.
The source $X(\omega)$ is easy to analyze:
it is a train of impulses,
where the spacing between impulses
characterizes the perceived pitch of the vowel.

The filter $H(\omega)$ has resonant frequencies known as \textbf{formants}.
These resonant frequencies create peaks in the graph of $H(\omega)$
at the position of each formant.
It is the shape of $H(\omega)$---or, equivalently,
the values of its formants---that gives a vowel sound its identity.
Importantly,

\begin{quote}
    \textbf{Claim.}
    A vowel's identity is \emph{not} determined by its \emph{exact, local} harmonic structure!
    Rather, it is determined by $H(\omega)$, which is represented by
    the ``spectral envelope'' of its \emph{rough, global} harmonic structure.
\end{quote}

You can play with the following (warning: AI-generated)
demo to see how formant values (F1, F2, etc.) affect vowel quality.

\begin{center}
\includewidget{lecture-16/speech-viz}
\end{center}

There are more than two formants,
but the first two are mostly sufficient
for uniquely defining vowel quality.
\begin{itemize}
    \item F1 is roughly higher when the tongue is lower,
    rather than higher.
    \item F2 is roughly lower when the tongue is retracted,
    rather than advanced.
\end{itemize}
In contrast, formants F4 and beyond
control subtler qualities such as vocal timbre.

An interesting corollary of all of this
is that higher-pitched vowels are harder to identify---they
give a much blurrier picture of the spectral envelope $H(\omega)$
because the train of impulses is spaced much further apart.
(Hear it for yourself in the demo!)

\begin{remark}
    When whispering, vocal cords are no longer in use,
    so $X(\omega)$ now represents a turbluent stream of air
    with no consistent, periodic harmonic structure at all.
    The result is still the same, though:
    the formants of $H(\omega)$ reveal themselves
    in the output $Y(\omega)$, determining vowel identity.
\end{remark}

\end{document}
